\chapter{Inline mode and the witness chain}
\label{ch:inline}

Most TUI libraries take over the entire terminal: they switch to the
\textbf{alternate screen}, paint, restore the original screen on
exit. \maya{} supports that --- it is the \texttt{Fullscreen}
rendering mode. But \maya{} also supports something rarer and far
more delicate: \textbf{inline mode}, where the UI lives \emph{at the
bottom of your normal scrollback}, redrawing in place but never
destroying what is above it.

Inline mode is what makes \agentty{} feel different from a
``classic'' TUI: when you type \code{agentty} in your shell, the chat
UI appears below your prompt without taking over the terminal, and
when it exits, your shell is exactly where you left it --- with the
conversation now permanent scrollback above the cursor. That
behaviour is, in principle, one wrong escape sequence away from
permanently corrupting the user's terminal history.

This chapter is a careful walk through how \maya{} guarantees it
\emph{doesn't}. We will look at the failure modes that real bugs in
the renderer produced (with citations), the four-part theorem
\maya{} commits to (\textbf{Scrollback Integrity}), and the five
layers of type-system machinery that close each leak in the theorem
into a compile error.

\section{The two modes, briefly}

\begin{lstlisting}
enum class Mode {
    Inline,      // raw mode, no alt screen, scrollback preserved
    Fullscreen,  // alt screen buffer, double-buffered cell diff
};
\end{lstlisting}

\textbf{Fullscreen mode} writes \code{ESC[?1049h} on startup
(``enter alt screen''), paints absolute-positioned cells, writes
\code{ESC[?1049l} on exit (``leave alt screen''). The user's
scrollback is bytes the renderer never touches. Easy.

\textbf{Inline mode} stays on the normal screen. The cursor starts
wherever the shell left it; the UI occupies some number of rows
\emph{below the cursor}; when the program quits, those rows remain
(or are erased and replaced). On every frame, the renderer must
re-paint the same region without ever writing a byte that affects
the rows above. The shell's history, your previous commands, the
output of \code{ls} you ran two minutes ago --- all immutable.

The naive question: what could go wrong? Two answers, both with
real bug reports in the repository.

\section{Failure mode 1: ghosting from over-tall paint}

The first real bug, documented in
\file{maya/docs/internals/inline-autogrow.md}, came from trying to
optimise the renderer. The legacy code pattern was: paint into the
existing canvas, measure how tall the content turned out, resize and
paint again if it overflowed. Two paint passes when content was big;
the first one always discarded.

The optimisation seemed obvious: ask the layout engine to compute
the natural height \emph{first} (with an effectively infinite height
constraint), resize, then paint once.

\begin{lstlisting}
constexpr int kBigH = 1 << 20;        // "infinite"
build_layout_tree(root, layout_nodes, theme);
layout::compute(layout_nodes, root_idx, w, kBigH);
const int natural_h = layout_nodes[root_idx].computed.size.height.raw();
if (canvas.height() < natural_h + 8) canvas.resize(w, natural_h + 8);
paint_element(root, canvas, pool, layout_nodes, root_idx, 0, 0);
\end{lstlisting}

One paint, correctly-sized canvas, no out-of-bounds window for the
component cache to capture from. Everything looks fine.

\textbf{It ghosted.} Interactive sessions started showing rows
duplicated and stacked into scrollback.

\subsection{Why}

The layout engine's \code{height} parameter is the \emph{available
height} for the root container. That value flows down to children
as their containing block. A child with \code{flex-grow > 0} or
\code{align-items: stretch} distributes a share of that available
height. Under \code{kBigH}, ``remaining height'' is pathologically
large: elastic children claim huge vertical extents that they would
never have claimed under a realistic constraint.

So the painted frame was a 50-row layout on a 24-row terminal.

\subsection{The downstream catastrophe}

After paint, \code{InlineFrameState::prev\_rows} records the painted
height (50). The next frame calls \code{compose\_inline\_frame},
which positions the cursor relative to where it left off:

\begin{lstlisting}
const int cursor_row_start = prev_rows - 1;            // 49
const int delta = first_changed - cursor_row_start;    // e.g. -45
if (delta < 0) {
    int up = std::min(-delta, prev_on_screen - 1);     // capped at 23
    if (up > 0) ansi::write_cursor_up(out, up);
    out += '\r';
}
\end{lstlisting}

\code{prev\_on\_screen = min(prev\_rows, term\_h) = 24}. The
\code{prev\_on\_screen - 1} cap is essential --- without it, the
cursor would walk \emph{up into scrollback-committed history},
overwriting the user's earlier output. With it, when we genuinely
need to move 45 rows up to reach \code{first\_changed}, we move
only 23. The cursor lands at the wrong physical line.

The per-row emit loop then proceeds as if the cursor were at
\code{first\_changed}:

\begin{lstlisting}
for (int y = first_changed; y <= last_row_to_visit; ++y) {
    if (y > first_changed) out += "\r\n";
    emit_row_diff(y, ...);
}
\end{lstlisting}

Every row gets written 22 lines too far down. Old content above
stays visible; new content overwrites whatever was previously at
those positions. When the content next shrinks back to a normal
size, the cycle repeats with a different misalignment --- and rows
that have already been pushed past the bottom into native
scrollback are gone, permanently.

\subsection{The fix}

Keep the two-compute pattern, but skip the wasted first paint:

\begin{lstlisting}
build_layout_tree(root, layout_nodes, theme);
layout::compute(layout_nodes, root_idx, w, canvas.height());   // realistic budget
const int needed_h = layout_nodes[root_idx].computed.size.height.raw();
if (canvas.height() < needed_h + 8) {
    canvas.resize(w, needed_h + 8);
    layout::compute(layout_nodes, root_idx, w, canvas.height());  // recompute
}
paint_element(root, canvas, pool, layout_nodes, root_idx, 0, 0);  // single paint
\end{lstlisting}

The recompute is what produces the painted layout, with the resized
height as the constraint --- byte-for-byte equivalent to legacy's
second \code{render\_tree} call, but with the first paint elided.

The doc concludes with the \textbf{invariant for future readers}:

\begin{quote}\itshape
The layout that gets painted must come from a \code{layout::compute}
call whose height parameter is at least as large as the canvas about
to receive the paint.
\end{quote}

That sentence is, in the larger sense, an unenforced human
discipline. A future contributor could ``simplify'' the code right
back into the broken shape, and CI would not catch it. We will
return to this kind of fragility in a moment --- it is why the
witness chain exists.

\section{Failure mode 2: \code{commit\_prefix} over-claim}

The second bug class, documented in \file{scrollback-corruption-audit.md},
is more insidious. It begins with a feature: a long-running chat UI
that streams turns should eventually \emph{commit} old turns to
scrollback, so that the live frame stays small and fresh content
keeps painting at the bottom.

The Cmd that does this is \code{Cmd::commit\_scrollback(int rows)}.
The handler, in the legacy implementation, simply decremented
\code{prev\_rows} and \code{memmove}-d \code{prev\_cells} up:

\begin{lstlisting}
// src/render/serialize.cpp (legacy)
prev_rows -= rows;
// (no bytes written, no escape emitted, no verification)
\end{lstlisting}

The implicit contract: ``\code{rows} is the count of rows that the
terminal \emph{already scrolled out of the viewport on its own},
because we emitted more \code{\textbackslash r\textbackslash n}s
than the terminal could fit.''

\subsection{Why a guessed row count corrupts everything}

\agentty{}'s \code{maybe\_virtualize} called

\begin{lstlisting}
return Cmd<Msg>::commit_scrollback(kSliceChunk * kRowsPerDroppedMessageLower);
\end{lstlisting}

The arithmetic was a \emph{guess} at the rendered height of dropped
messages, not a count of physical viewport scrolls. Concrete
scenario from the audit:

\begin{itemize}[leftmargin=*]
  \item 50-row terminal, 8 turns at $\sim$3 rows each:
    \code{prev\_rows $\approx$ 24}.
  \item All 24 rows are on screen; zero rows have scrolled.
  \item \code{commit\_inline\_prefix(8)} runs: \code{prev\_rows}
    drops from 24 to 16; \code{prev\_cells} is shifted up by 8.
  \item Physical viewport: unchanged. The renderer's model and
    reality now disagree by 8 rows.
\end{itemize}

Next compose:

\begin{itemize}[leftmargin=*]
  \item \code{cursor\_row\_start = 16 - 1 = 15}.
  \item Physical cursor: actually at row 23.
  \item Every emit lands 8 rows below where the renderer thinks.
  \item Later, when \code{\textbackslash r\textbackslash n} finally
    scrolls the viewport, the rows that scroll into native
    scrollback are \textbf{the stale ghost rows} --- the renderer's
    incorrect model, frozen into the user's permanent history.
\end{itemize}

That is the corruption. It is not recoverable: bytes already
committed to native scrollback are out of the renderer's reach.

\subsection{The structural fix}

Two paths, both shipped:

\begin{enumerate}[leftmargin=*]
  \item \textbf{Clamp inside the handler.} The Cmd handler now
    enforces \code{min(rows, max(0, prev\_rows - term\_h))}: a
    request to commit more rows than have actually overflowed the
    viewport becomes a no-op rather than a corruption.
  \item \textbf{Prefer the typed alternative.}
    \code{Cmd::commit\_scrollback\_overflow} takes no row count;
    it asks the renderer to commit precisely the rows that have
    provably overflowed (\code{prev\_rows - term\_h}). The caller
    cannot get the number wrong because there is no number.
\end{enumerate}

Path 1 makes the bug class harder to trigger; path 2 makes it
impossible. The library now prefers (2).

\section{The theorem: Scrollback Integrity}

After enough of these one-off fixes, the \maya{} authors did
something unusual: they wrote down the \emph{theorem} the renderer
should satisfy and then engineered the type system so violations
become compile errors.

\begin{theory}
\textbf{Scrollback Integrity.} For any well-typed \maya{} program
$P$, for any sequence of host events $e_1 \ldots e_n$, the bytes
the terminal emulator commits to its native scrollback are
\emph{exactly} the bytes \code{serialize} would produce for the
overflow-portion of past frames, in order. No row is ever targeted
by an escape sequence emitted by \maya{} once it has overflowed the
viewport.
\end{theory}

The theorem decomposes into four sub-claims, each closed by a
different layer of the chain:

\begin{itemize}[leftmargin=*]
  \item \textbf{T0 (canvas-cache integrity).} The
    \code{last\_content\_col(y)} and \code{max\_content\_row()}
    metadata a diff trusts is a true function of the cell buffer.
  \item \textbf{T1 (no-ghost).} No cursor-positioning escape
    targets a row above the viewport top.
  \item \textbf{T2 (no-corruption).} No cell-paint escape
    disagrees with \code{paint(view(model))} at the corresponding
    canvas coordinate.
  \item \textbf{T3 (monotone scrollback).} Once a row leaves the
    viewport via \code{\textbackslash r\textbackslash n} overflow,
    no subsequent render targets it.
\end{itemize}

T0 closes the ``stale tail'' ghost class (header shrinks; old tail
not cleared). T1 closes the ``cursor walks into history''
class. T2 closes the ``wire and shadow disagree'' class. T3 closes
the ``commit\_prefix over-claim'' class above.

Together: the bytes in emulator scrollback at any moment are the
union of \emph{truths committed so far} and \emph{truth currently on
screen}, both byte-identical to what the canvas claims. The
emulator's scrollback is what the model says it is. End of theorem.

\subsection{The honest hypotheses}

A theorem with no hypotheses is suspicious. \maya{}'s claim holds
\emph{only} when four assumptions hold:

\begin{itemize}[leftmargin=*]
  \item \textbf{A1 --- Terminal conformance.} The emulator honours
    ECMA-48 cursor control, DEC 2026 synchronised output if
    advertised, EL, ED, DECAWM. A buggy emulator that misinterprets
    \code{CSI A} is out of scope.
  \item \textbf{A2 --- Kernel write fidelity.} \code{write(2)}'s
    return value accurately reflects bytes delivered. POSIX
    guarantee; treat as axiom.
  \item \textbf{A3 --- Exclusive fd ownership.} \maya{}'s
    \code{Writer} is the only entity in the process writing to
    fd 1. Host code that calls \code{std::cout} or
    \code{fwrite(stdout)} outside \code{Cmd::print} or
    \code{Runtime::with\_external\_io} bypasses the entire chain.
    Cannot be enforced by the type system because \cpp{} does not
    let us seal \code{stdout}.
  \item \textbf{A4 --- No silent memory corruption.} Between
    \code{verify\_shadow} returning a witness and
    \code{compose\_inline\_frame\_v2} consuming it, the
    \code{prev\_cells} buffer is not mutated by use-after-free, an
    overflow in adjacent code, or a cosmic ray. An FNV-1a re-hash
    inside compose catches mutation with probability $1 - 2^{-64}$.
\end{itemize}

\section{The witness chain: five layers}

Each sub-claim is closed by a distinct piece of type machinery.
None of them depends on a comment, a convention, or a code review
catching it. The compiler is the proof assistant.

\subsection{Layer 0 --- \code{CanvasWitness} (T0)}

\textbf{Header:} \file{maya/include/maya/render/canvas\_witness.hpp}.

A \code{Canvas} caches metadata: per-row last-non-blank column
(\code{last\_content\_col\_[y]}) and overall max-content row
(\code{max\_y\_}). The diff loop reads these to skip blank tails ---
an enormous optimisation for typical UIs that are mostly empty
columns to the right of text.

The risk: if the cache lies (claims a longer tail than the cells
actually contain), the diff can re-emit cells that should be blank,
producing the classic ``stacked header, stale tail'' ghost row.

\code{CanvasWitness} is move-only, has a private constructor, and is
friended only to \code{verify\_canvas}. The verifier re-derives the
cache from the cells in a slow $O(W \cdot H)$ scan and returns
\code{nullopt} if they disagree. On match it records FNV-1a hashes
of both the cell buffer and the \code{(last\_col\_, max\_y\_)} tuple
and stamps the witness with them.

\begin{lstlisting}
class CanvasWitness {
    friend std::optional<CanvasWitness> verify_canvas(const Canvas&);
    Canvas* canvas_;
    uint64_t cells_hash_;
    uint64_t cache_hash_;
    CanvasWitness(Canvas* c, uint64_t ch, uint64_t kh)  // private
      : canvas_(c), cells_hash_(ch), cache_hash_(kh) {}
public:
    CanvasWitness(const CanvasWitness&) = delete;
    CanvasWitness(CanvasWitness&&) noexcept = default;
    uint64_t cells_hash() const noexcept { return cells_hash_; }
    uint64_t cache_hash() const noexcept { return cache_hash_; }
};
\end{lstlisting}

The \code{diff()} function that consumes the witness re-hashes the
moved-in canvas and \code{std::abort}s on mismatch. The renderer is
single-threaded; a mismatch between two consecutive calls signals
memory corruption, and there is no recovery path the renderer would
trust.

The host's only recovery if \code{verify\_canvas} returns
\code{nullopt} is \code{clear\_row()} plus full repaint --- never
``emit anyway.'' The escape hatch does not exist by construction.

\subsection{Layer 1 --- \code{WireRow} and \code{Viewport} (T1)}

\textbf{Header:} \file{maya/include/maya/render/wire\_row.hpp}.

Cursor coordinates are saturating. A \code{WireRow} is a wrapper
around an \code{int} whose constructor clamps:

\begin{lstlisting}
class WireRow {
    friend class Viewport;
    int row_;
    WireRow(int row, int top, int bottom)
      : row_(std::clamp(row, top, bottom - 1)) {}
public:
    int raw() const noexcept { return row_; }
};
\end{lstlisting}

The only public producers:

\begin{itemize}[leftmargin=*]
  \item \code{Viewport::for\_fresh\_frame(term\_h)} --- a fresh
    \code{Viewport} with the live region $[0,$ \code{term\_h}$)$.
  \item \code{Viewport::for\_redraw(term\_h, wire\_cursor\_rows)}
    --- the redraw variant, used after a frame has already started.
  \item \code{Viewport::row(int y)} --- the sole producer that takes
    an arbitrary absolute row index; clamps via \code{std::clamp}.
  \item \code{WireRow::up(n)} / \code{down(n)} --- arithmetic that
    saturates inside the same $[top, bottom)$.
\end{itemize}

\code{emit\_move(out, from, to)} is the sole emitter of cursor-up /
cursor-down escapes, and its parameters are typed \code{WireRow}.
There is no public \code{WireRow(int)} constructor, no
\code{data()}, no escape hatch.

\textbf{Proof of T1.} The set of cursor targets any \code{emit\_move}
can produce is precisely the set of \code{WireRow} values; the set
of \code{WireRow} values is precisely \code{Viewport::[top, bottom)}.
Targeting a row above the viewport top is, by construction,
unrepresentable.

This is the same flavour of compile-time proof as agentty's
permission policy from earlier: encode the rule as a type, make
violations uncompilable. The \code{kBigH} ghosting story from
earlier in this chapter could not happen against this layer ---
even if the renderer miscomputed the cursor delta, the resulting
\code{WireRow} would have saturated at the viewport top and no
escape would have aimed above it.

\subsection{Layer 2 --- \code{ContentRows} (closes one wrong-source class)}

\textbf{Header:} \file{maya/include/maya/render/serialize.hpp}.

\code{ContentRows} is an \code{int}-wrapper with a private
constructor. The sole producer is \code{content\_rows(canvas)},
which calls \code{content\_height(canvas)} on the canvas about to be
diffed and binds the result to that canvas's pointer.

\begin{lstlisting}
class ContentRows {
    friend ContentRows content_rows(const Canvas&);
    int rows_;
    const Canvas* src_;
    ContentRows(int r, const Canvas* s) : rows_(r), src_(s) {}
public:
    int raw()              const noexcept { return rows_; }
    const Canvas* source() const noexcept { return src_; }
};
\end{lstlisting}

\code{compose\_inline\_frame\_v2} consumes \code{ContentRows}. The
historical bug class --- computing the row count from the wrong
source (the layout's height, a hard-coded constant, a stale
\code{prev\_rows} from a different canvas, an agentty caller's
guess of how many rows are about to scroll) --- becomes
\emph{uncompilable}. There is exactly one way to construct a
\code{ContentRows}, and it forces the value to come from the canvas
the renderer is about to paint from.

\subsection{Layer 3 --- \code{ShadowWitness} (T2)}

\textbf{Header:} \file{maya/include/maya/render/serialize.hpp}.

The diff inside \code{compose\_inline\_frame\_v2} reads
\code{state.prev\_cells} --- the \emph{shadow buffer}, a mirror of
what the renderer believes is currently on the terminal --- and
emits ANSI for the cells that differ from the new canvas. If
\code{prev\_cells} disagrees with reality (use-after-free, an
unrelated bug overwriting it, an out-of-bounds write from another
subsystem), the diff emits the wrong bytes.

\code{ShadowWitness} is the move-only token that proves
\code{prev\_cells} matches the recorded shadow hash. The sole
producer is \code{verify\_shadow(state)}, which re-hashes
\code{state.prev\_cells} and returns \code{nullopt} on disagreement:

\begin{lstlisting}
class ShadowWitness {
    friend std::optional<ShadowWitness> verify_shadow(InlineFrameState&);
    InlineFrameState* state_;
    uint64_t hash_at_issue_;
    ShadowWitness(InlineFrameState* s, uint64_t h)         // private
      : state_(s), hash_at_issue_(h) {}
public:
    ShadowWitness(const ShadowWitness&) = delete;
    ShadowWitness(ShadowWitness&&) noexcept = default;
    uint64_t hash_at_issue() const noexcept { return hash_at_issue_; }
};
\end{lstlisting}

\code{compose\_inline\_frame\_v2} consumes a \code{ShadowWitness\&\&}.
Inside, it re-hashes the moved-in state and \code{std::abort}s on
mismatch. The window between \code{verify\_shadow} returning a
witness and compose re-checking is small --- both calls happen on
the same thread, in the same function --- so under hypothesis A4
(no silent memory corruption) the cells the diff reads are the cells
the terminal is showing, and the emitted bytes are the byte-for-byte
difference.

\subsection{Layer 4 --- \code{FrameBytes} (atomic state advance)}

\textbf{Header:} \file{maya/include/maya/render/frame\_bytes.hpp}.

\code{compose\_inline\_frame\_v2} returns a \code{FrameBytes} ---
the bytes the frame would emit, plus the \code{InlineFrameState}
the caller should hold \emph{if the bytes ship}. It is move-only.
The only consumers:

\begin{itemize}[leftmargin=*]
  \item \code{commit\_to(writer) \&\& -> CommitOutcome} --- a
    \code{std::variant} of \code{commit::Synced\{state\}} (full
    delivery) or \code{commit::Stale\{state, reason\}} (hard I/O
    error).
  \item \code{abandon() \&\& -> commit::Stale\{state, ok()\}} ---
    explicit drop.
\end{itemize}

Because the return type is a \code{std::variant} and structured
bindings on \code{variant} are ill-formed, every call site must
\code{std::visit} both arms. \textbf{There is no success-only code
path.}

\begin{lstlisting}
auto bytes = std::move(synced).render(canvas, content_rows(canvas),
                                       term_h, pool, *std::move(wit));
auto outcome = std::move(bytes).commit_to(*writer);
return std::visit(overload{
    [](commit::Synced&& s) { return RenderResult{std::move(s.state)}; },
    [](commit::Stale&&  s) { return RenderResult{std::move(s.state),
                                                 s.reason}; }
}, std::move(outcome));
\end{lstlisting}

The historical bug class --- ``state advanced; \code{write} returned
partial; next compose diffed against a state the wire doesn't
show'' --- becomes a compile error, because the only way to obtain a
\code{Synced} state is the return value of a successful
\code{commit\_to}, and the only way to obtain a \code{Stale} state
is from either \code{commit\_to}'s I/O-error arm or
\code{abandon()}. State advance and byte delivery are fused into
one operation whose outcome is a value the caller must destructure.

\subsection{Layer 5 --- \code{InlineFrame<Tag>} (linear chain)}

\textbf{Header:} \file{maya/include/maya/render/inline\_frame.hpp}.

The apex. \code{InlineFrame<Tag>} is a phantom-tagged type-state.
Five tags:

\begin{center}
\small
\begin{tabular}{ll}
\textbf{Tag} & \textbf{Meaning} \\
\hline
\code{Empty}  & no frame yet \\
\code{Fresh}  & first frame just rendered \\
\code{Synced} & wire and shadow agree (this is the steady state) \\
\code{Stale}  & wire and shadow may disagree (after an I/O error) \\
\code{Sealed} & program is exiting; no further render \\
\end{tabular}
\end{center}

Each tag is a distinct class (template specialisation); each is
move-only; each has a private constructor friended only to its
predecessors in the transition graph:

\begin{center}
\small
\begin{tabular}{lll}
\textbf{From} & \textbf{Op} & \textbf{To} \\
\hline
\code{Empty}    & \code{seed()}                  & \code{Fresh}                          \\
\code{Fresh}    & \code{render(...)}             & \code{Synced} or \code{Stale}         \\
\code{Synced}   & \code{verify()}                & \code{optional<ShadowWitness>}        \\
\code{Synced}   & \code{render(..., wit\&\&)}     & \code{Synced} or \code{Stale}         \\
\code{Synced}   & \code{demote\_to\_stale()}      & \code{Stale}                          \\
\code{Synced}   & \code{commit(marker)}          & \code{Synced} (scrollback advance)    \\
\code{Stale}    & \code{render(...)}             & \code{Synced} or \code{Stale}         \\
\code{any}      & \code{finalize(out)}           & \code{Sealed}                         \\
\end{tabular}
\end{center}

Compile-time consequences:

\begin{itemize}[leftmargin=*]
  \item \code{InlineFrame<Synced>::render} requires
    \code{ShadowWitness\&\&}. You cannot render a Synced frame
    without first calling \code{verify()}.
  \item \code{InlineFrame<Sealed>} has no \code{render} method ---
    rendering after \code{finalize} is a compile error.
  \item \code{InlineFrame<Empty>::render} does not exist ---
    rendering before \code{seed} is a compile error.
  \item All five specialisations are non-copyable ---
    holding two parallel views of the state is a compile error.
\end{itemize}

The non-copyability is enforced by \code{static\_assert}:

\begin{lstlisting}
static_assert(!std::is_copy_constructible_v<InlineFrame<Empty>>);
static_assert(!std::is_copy_constructible_v<InlineFrame<Fresh>>);
static_assert(!std::is_copy_constructible_v<InlineFrame<Synced>>);
static_assert(!std::is_copy_constructible_v<InlineFrame<Stale>>);
static_assert(!std::is_copy_constructible_v<InlineFrame<Sealed>>);
static_assert(!std::is_copy_constructible_v<ShadowWitness>);
static_assert(!std::is_copy_constructible_v<FrameBytes>);
\end{lstlisting}

If a future refactor accidentally adds \code{= default} after
removing the move-only \code{delete}s, these \code{static\_assert}s
fire and the build breaks. The discipline is checked by the
compiler on every commit.

\subsection{Scrollback commit, formalised (T3)}

The Cmd that committed the legacy bug --- ``commit N rows to
scrollback'' --- has been completely re-typed. The new flow:

\begin{lstlisting}
auto marker = synced.scrollback_marker(rows);   // rows clamped to [0, prev_rows]
auto next   = std::move(synced).commit(std::move(marker));
\end{lstlisting}

\code{ScrollbackMarker} is move-only and consumed exactly once. The
returned \code{InlineFrame<Synced>} has \code{prev\_rows} strictly
less than or equal to the prior. The set of row coordinates any
future cursor move can target is bounded above by
\code{prev\_rows\_at\_emit + W} where \code{W} is the wire-cursor
row count from the most recent successful emit.

T3 follows: once a row leaves the viewport via overflow, no future
emit can target it, because no \code{Viewport} produced from any
later state contains it, and \code{WireRow} closes the cursor
inside the viewport.

\subsection{What the chain does \emph{not} claim}

Honest scope:

\begin{itemize}[leftmargin=*]
  \item It does not prevent bugs \emph{inside} the diff algorithm.
    If a per-row emit miscalculates a span or wide-char snap fails
    on a new emoji range, the wrong bytes go out for the right
    state. Type-state cannot witness algorithmic correctness;
    property-based fuzzing of the diff is the right complement.
  \item It does not prevent the terminal emulator from dropping
    bytes (A1).
  \item It does not prevent host code from writing to fd 1 directly
    (A3).
  \item It does not prevent silent memory corruption (A4); the
    FNV-1a re-checks are the probabilistic defense.
\end{itemize}

These are not handwaves. Each is a hypothesis named in the proof.

\section{Where the code lives}

A quick map for when you want to read the source:

\begin{center}
\small
\begin{tabular}{ll}
\textbf{File} & \textbf{Role} \\
\hline
\file{maya/include/maya/render/canvas\_witness.hpp} & Layer 0 --- T0 \\
\file{maya/include/maya/render/wire\_row.hpp}        & Layer 1 --- T1 \\
\file{maya/include/maya/render/serialize.hpp}        & Layers 2--3 --- ContentRows, ShadowWitness \\
\file{maya/include/maya/render/frame\_bytes.hpp}     & Layer 4 --- atomic state advance \\
\file{maya/include/maya/render/inline\_frame.hpp}    & Layer 5 --- linear chain \\
\file{maya/src/render/serialize.cpp}                  & \code{compose\_inline\_frame\_v2} \\
\file{maya/src/render/inline\_frame.cpp}              & runtime transitions \\
\file{maya/docs/internals/witness-chain.md}           & this chapter's source material \\
\file{maya/docs/internals/scrollback-corruption-audit.md} & every prior bug, with citations \\
\file{maya/docs/internals/inline-autogrow.md}         & the ghosting trap, in detail \\
\file{maya/tests/test\_wire\_row.cpp}                 & 10 tests + saturation proofs \\
\file{maya/tests/test\_inline\_frame.cpp}             & 5 happy paths + \code{static\_assert}s \\
\end{tabular}
\end{center}

\section{Why all this matters}

It is reasonable to ask: this is a lot of machinery for what is, in
the end, ``don't write the wrong escape sequence.'' Why?

The answer is the same one you find in every chapter of this book.
\maya{} is a library; libraries are used by people who did not
write them; the bugs that matter most are the ones the user never
triggers in dev because their terminal happens to be 80x24 and
their chat fits in a viewport. They appear in production --- on a
50-line terminal, after a long session, after a virtualising
\code{Cmd::commit\_scrollback}, in a way that destroys the user's
permanent terminal history.

The fixes you'd write to those bugs as one-off patches eventually
diverge from the discipline they enforce. A new contributor will
read the patched code, not the patch, and re-introduce the bug. The
witness chain makes the discipline a \emph{property of the type
system}. The compiler enforces it on every build, in every
contributor's editor, without anyone having to remember.

If you take one thing from this chapter, take that:
\textbf{discipline encoded in types survives refactoring;
discipline written in prose does not}.

\begin{recap}
Inline mode is rendering into the user's normal scrollback without
ever destroying what is above the UI. Two historical bug classes
--- \emph{ghosting from over-tall paint} and \emph{over-claimed
scrollback commit} --- showed how brittle the convention-only
approach was. \maya{}'s answer is the \textbf{witness chain}: a
four-part theorem (\emph{Scrollback Integrity}) decomposed into T0
(canvas cache truth), T1 (no cursor escape above viewport), T2
(no cell paint disagreeing with the model), T3 (monotone
scrollback), each closed by a move-only typed token (\code{CanvasWitness},
\code{WireRow}, \code{ShadowWitness}, \code{FrameBytes}) and
ultimately by the \code{InlineFrame<Tag>} linear chain. Every
prior bug class is now a compile error.
\end{recap}

\section*{Workshop}

\begin{enumerate}[leftmargin=*]
  \item Read \file{maya/include/maya/render/wire\_row.hpp}. Find
    the private constructor, the \code{std::clamp} call, and the
    list of friends. Try to construct a \code{WireRow} outside
    \code{Viewport}; convince yourself the compiler refuses.
  \item Read \file{maya/docs/internals/witness-chain.md}. Identify
    which paragraph of the theorem corresponds to T0, T1, T2, T3.
    Match each to its layer in the chain.
  \item Read \file{maya/docs/internals/scrollback-corruption-audit.md},
    Finding \#1 (\code{commit\_prefix} over-claim). Match the
    failure mode in the audit to the layer in the chain that now
    prevents it.
  \item Look at \file{maya/tests/test\_inline\_frame.cpp}. Identify
    the \code{static\_assert}s that pin non-copyability. Imagine
    deleting one of the \code{delete}d copy constructors; predict
    which \code{static\_assert} fires.
\end{enumerate}
